

\section{Problem Formulation}
\label{sec:tkg}

We first describe notations for building our model and problem definition, and then we define the joint distribution of temporal events. 
% In the next section, we will describe details of \method, our autoregressive structure inference framework. 

\para{Notations and Problem Definition.}
We consider a \textit{temporal knowledge graph} as a multi-relational, directed graph with time-stamped edges between nodes (entities).
An \textit{event} is defined as a time-stamped edge, \textit{i.e.}, (subject entity, relation, object entity, time) and is denoted by a quadruple $(\bs,\br,\bo,t)$ or $(\bs_t, \br_t, \bo_t)$. We denote a set of events at time $t$ as ${\calE}_t$. 
In our setup, the timestamps are discrete integers and used for the relative order of graphs or events.
A TKG is built upon \textit{a sequence of event quadruples} ordered ascending based on their timestamps, \textit{i.e.}, $\{\calE_{\bt}\}_t = \{(\bs_i,\br_i,\bo_i,t_i)\}_i$ ($t_i < t_j, \forall i<j$), 
where each time-stamped edge has a direction pointing from the subject entity to the object entity.\footnote{The same triple $(\bs, \br, \bo)$ may occur multiple times in different timestamps, yielding different event quadruples.}  
The goal of \textit{learning generative models of events} is to learn a distribution $\prob(\calE)$ over TKGs, based on a set of observed event sets $\{\calE_1, ..., \calE_t\}$. 



\para{Approach Overview.}
The key idea of our approach is to learn temporal dependency from the sequence of graphs and local structural dependency from the neighborhood (Fig.~\ref{fig:renet}). 
Formally, we represent TKGs as sequences, and then build an autoregressive generative model on the sequences.
To this end, \method defines the joint probability of concurrent events (or a graph), which is conditioned on all the previous events.
Specifically, \method consists of a Recurrent Neural Network (RNN) as a recurrent event encoding module and a neighborhood aggregation module to capture the information of graph structures.
We first start with the definition of joint distribution of temporal events.



\para{Modeling Joint Distribution of Temporal Events.}
We define the joint distribution of all the events $\calE = \{\calE_{1}, ..., \calE_{T}\}$ in an autoregressive manner.
% , i.e., $\prob(\calE) = \prod_{t=1}^T{\prob(\calE_t|\calE_{t-m:t-1})}$. 
Basically, we decompose the joint distribution into a sequence of conditional distributions, $\prob(\calE_t|\calE_{t-m:t-1})$), where 
we assume the probability of the events at a time step, $\calE_t$, depends on the events at the previous $m$ steps, $\calE_{t-m:t-1}$. 
For each conditional distribution $\prob(\calE_t|\calE_{t-m:t-1})$, we further assume that the events in $\calE_t$ are mutually independent given the previous events $\calE_{t-m:t-1}$. In this way, the joint distribution can be rewritten as follows:

\small
\begin{multline}
\label{eq:prob_event_t}
    \prob(\calE)= \prod_{t}\prod_{(\bs_t,\br_t,\bo_t)\in\calE_{t}}  \prob(\bs_t,\br_t,\bo_t|\calE_{t-m:t-1})\\
    = \prod_t \prod_{(\bs_t,\br_t,\bo_t)\in\calE_t} 
    \prob(\bs_t |\calE_{t-m:t-1})
    \cdot \prob(\br_t |\bs_t, \calE_{t-m:t-1})\\
    \cdot \prob(\bo_t|\bs_t,\br_t,\calE_{t-m:t-1}).
\end{multline}
\normalsize

% \begin{equation}
% \label{eq:prob_event_t}
%     \small \prob(\calE)= \prod_{t}\prod_{(\bs_t,\br_t,\bo_t)\in\calE_{t}}  \prob(\bs_t,\br_t,\bo_t|\calE_{t-m:t-1}).
% \end{equation}


% We factorize the joint probability of a triplet into three conditional distribution. 
From these probabilities, we generate triplets as follows.
Given all the past events $\calE_{t-m:t-1}$, we first sample a subject entity $\bs_t$ through $\prob(\bs_t |\calE_{t-m:t-1})$. Then we generate a relation $\br_t$ with $\prob(\br_t |\bs_t, \calE_{t-m:t-1})$, and finally the object entity $\bo_t$ is generated by $\prob(\bo_t|\bs_t,\br_t,\calE_{t-m:t-1})$.\footnote{We can also first sample an object entity in this process. Details are omitted for brevity.}
% Note that the probability is factorized in the order of $\bo_t, \br_t, \bs_t$. 

% We omit it for brevity.


% In this work,
% we assume that $\prob(\bo_t|\allowbreak\bs_t,\br_t,\calE_{t-m:t-1})$ and $\prob(\br_t | \allowbreak \bs_t, \calE_{t-m:t-1})$ depend only on events that are related to $\bs$, and focus on modeling the following joint probability:
% \begin{multline}
% \label{eq:prob_factor}
%     \prob(\bs_t,\br_t,\bo_t|\calE_{t-m:t-1}) \\=  \prob(\bs_{t}|\calE_{t-m:t-1})
%     \cdot \prob(\br_{t}|\bs,\teN^{(\bs)}_{t-m:t-1})
%     \cdot \prob(\bo_{t}|\bs,\br,\teN^{(\bs)}_{t-m:t-1}),
% \end{multline}
% where $\calE_{t}$ becomes $\teN^{(\bs)}_t$ which is a set of neighboring entities interacted with subject entity $\bs$ under \textit{all} relations at time $t$. 
% For $\prob(\bs_{t}|\calE_{t-m:t-1})$, all event sets should be considered since subject $\bs$ is not given. 
Next, we introduce how these probabilities are defined and parameterized in our method.


\section{Recurrent Event Network}
\label{sec:re-net}
In this section, we introduce our proposed method, Recurrent Event Network (\method).
\method consists of a Recurrent Neural Network (RNN) as a recurrent event encoder (Sec.~\ref{sec:encoder}) for temporal dependency and a neighborhood aggregator (Sec.~\ref{sec:aggre}) for graph structural dependency.
We also discuss parameter learning of \method and define multi-step inference for distant future by sampling intermediate graphs in a sequential manner (Sec.~\ref{sec:learn}). 

% \xiang{move and fill up some sentences to transit from previous section to here, and overview the subsections in this part.}

\subsection{Recurrent Event Encoder}
\label{sec:encoder}

To parameterize the probability for each event, \method introduces a set of \textit{global} representations as well as \textit{local} representations. The global representation $\bH_{t}$ summarizes the global information from the entire graph until time stamp $t$, which reflects the global preference on the upcoming events. In contrast, the local representations focus more on each subject entity $\bs$ or each pair of subject entity and relation $(\bs,\br)$, which capture the knowledge specifically related to those entities and relations. 
We denote the above local representations as $\bh_{t}(\bs)$ and $\bh_{t}(\bs,\br)$, respectively. 
The global and local representations capture different aspects of knowledge from the knowledge graph, which are naturally complementary, allowing us to model the generative process of the graph in a more effective way.

Based on the above representations, \method parameterizes $\prob(\bo_t| \allowbreak \bs,\br,\calE_{t-m:t-1})$ in the following way:
\begin{equation}
% \setlength{\belowdisplayskip}{0.5pt}
\label{eq:prob}
	%\prob_{\bo|\mathbf{s,r,\calE}}( \bo_{t}|\bs,\br,\teN^{(\bs)}_{t-m:t-1})=f_1(\es,\er,\bh_{t-1}(\bs, \br))[\bo_{t}],
	\small \prob( \bo_{t}|\bs,\br,\calE_{t-m:t-1}) \propto \exp \left([\es:\er:\bh_{t-1}(\bs, \br)]^\top \cdot  \bw_{\bo_{t}}\right),
\end{equation}
where $\es, \er \in \mbR^d$ are learnable embedding vectors specified for subject entity $\bs$ and relation $\br$. 
$\bh_{t-1}(\bs,\br)\in \mbR^d$ is the local representation for $(\bs,\br)$ obtained at time stamp $(t-1)$. Intuitively, $\es$ and $\er$ can be understood as static embedding vectors for subject entity $\bs$ and relation $\br$, whereas $\bh_{t-1}(\bs,\br)$ is dynamically updated at each time stamp. By concatenating both the static and dynamic representations, \method can effectively capture the semantic of $(\bs,\br)$ up to time stamp $(t-1)$. Based on that, we further compute the probability of different object entities $\bo_{t}$ by passing the encoding into our multi-layer perceptron (MLP) decoder.
We define the MLP decoder as a linear softmax classifier parameterized by $\{\bw_{\bo_{t}}\}$.



Similarly, we define probabilities for relations and subjects as follows:
% \meng{It is the first time that $H$ appear in the paper. We might need to add more explanation. Also, for all the vectors, we may use mathbf to represent them.}

\small
\begin{align}
\label{eq:prob2}
	\prob(\br_{t}|\bs,\calE_{t-m:t-1})&\propto \exp \left([\es:\bh_{t-1}(\bs))]^\top \cdot \bw_{\br_{t}} \right),\\
	\label{eq:prob3}
	\prob(\bs_{t}|\calE_{t-m:t-1})& \propto \exp \left( \bH_{t-1}^\top \cdot \bw_{\bs_{t}} \right),
\end{align}
\normalsize
where $\bh_{t-1}(\bs)$ focuses on the local information about $\bs$ in the past, and $\bH_{t-1} \in \mbR^d$ is a vector representation to encode global graph structures $\calE_{t-1:t-m}$. To predict what relations a subject entity will interact with  $\prob(\br_{t}|\bs,\calE_{t-m:t-1})$, we treat the static representation $\es$ as well as the dynamic representation $\bh_{t-1}(\bs)$ as features, and feed them into a multi-layer perceptron (MLP) decoder parameterized by $\bw_{\br_{t}}$. Besides, to predict the distribution of subject entities at time stamp $t$ (i.e., $\prob(\bs_{t}|\calE_{t-m:t-1})$), we treat the global representation $\bH_{t-1}$ as a feature, as it summarizes the global information from all the past graphs until time stamp $t-1$, which reflects the global preference on the upcoming events at time stamp $t$.

% Both the local representations $\bh_{t-1}(\bs, \br), \bh_{t-1}(\bs)$ and the global representation $\bH_{t-1}$ are dynamic, meaning that we will update them according to the observed events at each time stamp. 

The global representation $\bH_t$ is expected to preserve the global information about all the graphs up to time stamp $t$.
The local representations $\bh_t(\bs, \br)$ and $\bh_t(\bs)$ emphasize more on the local events related to each entity and relation.
Thus we define them as follows:
% We define the global representation $\bH_t$ and local representations $\bh_{t}(\bs,\br)$ and $\bh_{t}(\bs)$ as follows:
% \begin{equation}
% 	\label{eq:globalh}
% 	\small \bH_{t} = \textrm{RNN}^1(g(\calE_t), \bH_{t-1}).
% \end{equation}

\small
\begin{align}
    \bH_{t} &= \textrm{RNN}^1(g(\calE_t), \bH_{t-1}), \label{eq:globalh}\\
	\bh_{t}(\bs,\br) &= \textrm{RNN}^2(g(\teN^{(\bs)}_{t}), \bH_{t}, \bh_{t-1}(\bs,\br)),\\
	\bh_{t}(\bs) &= \textrm{RNN}^3(g(\teN^{(\bs)}_{t}), \bH_{t},\bh_{t-1}(\bs)),
\end{align}
\normalsize
where $g$ is an aggregate function which will be discussed in Section~\ref{sec:aggre} and $\teN^{(\bs)}_{t}$ stands for all the events related to $\bs$ at the current time step $t$.
We leverage a recurrent model $\textrm{RNN}$~\cite{Cho2014LearningPR} to update them.
% \footnote{Details are described in Section~\ref{append:eventseq} of appendix.}
The global representation takes the global graph structure $g(\calE_t)$ as an input.
$g(\calE_t)$ is an aggregation over all the events $\calE_t$ at time $t$. 
We define {\small$g(\calE_t) = \max(\{g(\teN_t^{(s)})\}_s)$}, which is an element-wise max-pooling operation over all {\small$g(\teN^{(\bs)}_{t})$}. The $g(\teN^{(\bs)}_{t})$ captures the local graph structure for subject entity $s$. 
The local representations are different from the global representations in two ways.
First, the local representations focus more on each entity and relation, and hence we aggregate information from events $\teN^{(\bs)}_{t}$ that are related to the entity. Second, to allow \method to better characterize the relationships between different entities, we treat the global representation $\bH_{t}$ as an extra feature in the definition, which acts as a bridge to connect different entities.

% The global representation $\bH_t$ is expected to preserve the global information about all the graphs up to time stamp $t$, so we update it as follows:
% % For the global representation $\bH_t$, it is expected to preserve the global information about all the graphs up to time stamp $t$, and therefore we update $\bH_t$ in the following recurrent way:
% \begin{equation}
% 	\label{eq:globalh}
% 	\small \bH_{t} = \textrm{RNN}^1(g(\calE_t), \bH_{t-1}).
% \end{equation}


% We leverage a recurrent model $\textrm{RNN}^1$~\cite{Cho2014LearningPR} to update $\bH_t$ based on the global representation $\bH_{t-1}$ at the previous time stamp as well as the information $g(\calE_t)$ from all the events at time stamp $t$. 
% Here, $g(\calE_t)$ is an aggregation function over all the events $\calE_t$ at time $t$. 
% More specifically, we define {\small$g(\calE_t) = \max(\{g(\teN_t^{(s)})\}_s)$}, which is an element-wise max-pooling operation over all {\small$g(\teN^{(\bs)}_{t})$}, where $g(\teN^{(\bs)}_{t})$ is the local information for each subject entity $s$. 
% We will explain the details shortly. In this way, $\bH_t$ is able to preserve both the past information and the information from the current entire events.

% For the local representations $\bh_t(\bs, \br), \bh_t(\bs)$, they emphasize more on the local events related to each entity and relation. Therefore, we use the following recurrent models for updating the local representations:

% \small
% \begin{align}
% 	\bh_{t}(\bs,\br) &= \textrm{RNN}^2(g(\teN^{(\bs)}_{t}), \bH_{t}, \bh_{t-1}(\bs,\br)),\\
% 	\bh_{t}(\bs) &= \textrm{RNN}^3(g(\teN^{(\bs)}_{t}), \bH_{t},\bh_{t-1}(\bs)),
% \end{align}
% \normalsize
% where $g$ is an aggregation function, and $\teN^{(\bs)}_{t}$ stands for all the events related to $\bs$ at the current time step $t$. $\textrm{RNN}^2$ and $\textrm{RNN}^3$ are two recurrent models based on Gated Recurrent Units~\cite{Cho2014LearningPR}, and details are described in Section~\ref{append:eventseq}.  The update rules are similar to those of the global representations but with two differences. First, the local representations focus more on each entity and relation, and hence we aggregate information from events $\teN^{(\bs)}_{t}$ that are related to the entity. Second, to allow \method to better characterize the relationships between different entities, we treat the global representation $\bH_{t}$ as an extra feature in the update rule, which acts as a bridge to connect different entities. With such update rules, the local representations can effectively encode the meaning and relationship of each entity and relation.

% The key of the above update rules is to design 
In the next section, we introduce how we design $g$ in \method.
% The update rules require a proper aggregation function $g$ to effectively encode events. However, it is nontrivial because for each subject entity $\bs$, it can interact with multiple relations and object entities at each time step $t$, \ie, the set {\small$\teN^{(\bs)}_{t}$} can contain multiple events. Next, we introduce how we design $g$ in \method.

\begin{figure}[tb!]
    \centering
    \includegraphics[width=1\columnwidth]{figures/compar_aggre.pdf}
    \caption{\textbf{Comparison of neighborhood aggregators.}
    The blue node corresponds to node $\bs$, red nodes are 1-hop neighbors, and green nodes are 2-hop neighbors. Different colored edges are different relations.
    Mean and attentive pooling aggregators do not differentiate different relations and do not encode 2-hop neighbors, whereas RGCN aggregator can incorporate information from multi-relational and multi-hop neighbors.
    }
    \label{fig:rgcnaggre}
\end{figure}



\subsection{Neighborhood Aggregators}
\label{sec:aggre}
% \meng{It might be better if we can explain the high-level intuition here. Otherwise, readers may not fully understand the motivation of the RGCN aggregator.}
% Recall that \method relies on an aggregation function $g$, which aims at summarizing the local information and the global information. 
In this section, we first introduce two simple aggregation functions: a mean pooling aggregator and an attentive pooling aggregator. These two simple aggregators only collect neighboring entities under the same relation $\br$. Then we introduce a more powerful aggregation function: a multi-relational aggregator.
We depict comparison on aggregators in Fig.~\ref{fig:rgcnaggre}.
% \textcolor{red}{Yu: IMO no need to italicize i.e. \& e.g. as they are common latin abbrv in english}
%Before going into a multi-relational graph aggregator, we first introduce a simple mean. 
% We assume that the next set of objects can be predicted with the previous object history under the same relation. \meng{I don't quite understand the sentence.}



\para{Mean Pooling Aggregator.}
The baseline aggregator simply takes the element-wise mean of representations in {\small $\{\eo: \bo \in \teN_t^{(\bs,\br)}\}$}, where {\small$\teN_t^{(\bs,\br)}$} is the set of objects that interacted with $\bs$ under $\br$ at $t$. 
% \textcolor{red}{Yu: Minor grammar mistake - "is to" is somehow improper}
But the mean aggregator treats all neighboring objects equally, and thus ignores the different importance of each neighbor entity. 

\para{Attentive Pooling Aggregator.}
We define an attentive aggregator based on the additive attention introduced in \citep{Bahdanau2015NeuralMT} to distinguish the important entities for $(\bs,\br)$. 
The aggregate function is defined as {\small $g(\teN_t^{(\bs,\br)}) = \sum_{\bo \in \teN_t^{(\bs, \br)}} \alpha_\bo \eo$},
where {\small $\alpha_\bo = \textrm{softmax}(\bv^\top \tanh(\bW(\es; \er; \eo)))$}.
{\small  $\bv \in \mathbb{R}^d$} and  {\small $\bW \in \mathbb{R}^{d \times 3d}$} are trainable weight matrices. 
By adding the attention function of the subject and the relation, the weight can determine how relevant each object entity is to the subject and the relation.


\para{Multi-Relational Graph (RGCN) Aggregator.} 
% \textcolor{red}{Yu: Is this the identical aggregator appeared in R-GCN paper?}
We introduce a multi-relational graph aggregator from \citep{Schlichtkrull2018ModelingRD}. 
This is a general aggregator that can incorporate information from multi-relational and multi-hop neighbors. Formally, the aggregator is defined as follows:
\begin{equation}
\label{eq:rgcn}
    \small g(\teN^{(\bs)}_t)= \bbh_s^{(l+1)} = \sigma\Big(\sum_{\br\in R}\sum_{\bo \in \teN_t^{(\bs,\br)}} \frac{1}{c_s} \bW_r^{(l)} \bbh_\bo^{(l)} + \bW_0^{(l)} \bbh_s^{(l)}\Big), 
\end{equation}
where initial hidden representations for each node ($\bbh_o^{(0)}$) are set to trainable embedding vectors ($\eo$) for each node and $c_s$ is a normalizing factor.
Details are described in Section~\ref{sec:rgcn} of appendix.

% \wj{I moved details to appendix.}

\begin{algorithm}[t]
\begin{small}
    \caption{\small Learning algorithm of \method}\label{alg:renet}
	\KwInput{Observed graph sequence: $\{\calE_{1}, ..., \calE_{t}\}$, Number of events to sample at each step: $M$.}
	\KwOutput{An estimation of the conditional distribution $\prob(\calE_{t+\Delta t}|\calE_{:t})$.}
	$t' \leftarrow t+1$\\
	\While{$t' \leq t+ \Delta t$}{
    Sample $M$ number of $\bs ~\sim \prob(\bs|\hat{\calE}_{t+1:t'-1},\calE_{:t})$ by equation \eqref{eq:prob3}.\label{alg:samples}\\
    Pick top-$k$ triples $\{(\bs_1,\br_1,\bo_1, t'),...,(\bs_k,\br_k,\bo_k, t')\}$ ranked by $\prob(\bs,\br,\bo|\hat{\calE}_{t+1:t'-1},\calE_{:t})$. \label{alg:sampletriple}\\
    ${\hat{G}}_{t'} \leftarrow \{(\bs_1,\br_1,\bo_1, t'),...,(\bs_k,\br_k,\bo_k, t')\}$ \label{alg:graph}\\
	$t' \leftarrow t'+ 1$
	}
	Estimate the probability of each event $\prob(\bs,\br,\bo|\hat{\calE}_{t+1:t+\Delta t-1},\calE_{:t})$.\\
	Estimate the joint distribution of all events $\prob(\calE_{t+\Delta t}|\hat{\calE}_{t+1:t+\Delta t-1},\calE_{:t})$ by equation \eqref{eq:prob_event_t}.
	\\
	\Return $\prob(\calE_{t+\Delta t}|\hat{\calE}_{t+1:t+\Delta t-1},\calE_{:t})$ as an estimation.
	\end{small}
\end{algorithm}


\subsection{Parameter Learning and Inference}
\label{sec:learn}
In this section, we discuss how \method is trained and infers events over multiple time stamps.

\para{Parameter Learning via Event Predictions.}
An object entity prediction given $(\bs,\br)$ can be viewed as a multi-class classification task, where each class corresponds to each object entity.
Similarly, relation prediction given $\bs$ and subject entity prediction can be considered as a multi-class classification task. 
Here we omit the notations for preceding events for brevity.
% To learn parameters and representations of entities and relations, we adopt a multi-class cross-entropy loss to the model's output.
Thus, the loss function is as follows:
% The loss function is comprised of three losses and is defined as:
\begin{equation}
    \small \mathcal{L} = -\sum_{(\bs,\br,\bo,t) \in \calE}  \log\prob(\bo_t|\bs_t,\br_t) +\lambda_1 \log\prob(\br_t|\bs_t) + \lambda_2 \log\prob(\bs_t),
    \label{loss}
\end{equation}
where $\calE$ is set of events, and $\lambda_1$ and $\lambda_2$ are importance parameters that control the importance of each loss term. $\lambda_1$ and $\lambda_2$ can be chosen depending on the task. 
If a task aims to predict  $o$ given $(s,r)$, then we can give small values to $\lambda_1$ and $\lambda_2$.
% Each probability in equation~\eqref{loss} is defined in equations \eqref{eq:prob}, \eqref{eq:prob2}, and \eqref{eq:prob3}, respectively.
% We apply teacher forcing for model training over historical data, i.e., we use the ground truth rather than the model's own prediction as the input of the next time step during training.






\para{Multi-step Inference over Time.} 
\method seeks to predict the forthcoming events based on the previous observations. 
Suppose that the current time is $t$ and we aim to predict events at time $t + \Delta t$ where $\Delta t>0$. 
Then the problem of multi-step inference can be formalized as inferring the conditional probability {\small $\prob (\calE_{t+\Delta t}|\calE_{:t})$}. 
The problem is nontrivial as we need to integrate over all $\calE_{t+1:t+\Delta t-1}$. To achieve efficient inference, we draw a sample of $\calE_{t+1:t+\Delta t-1}$, and estimate the conditional probability as follows:

\small
\begin{align*}
	&\prob (\calE_{t+\Delta t}|\calE_{:t})\nonumber\\
	&= \sum_{\calE_{t+1:t+\Delta t-1}} \prob (\calE_{t+\Delta t}, \calE_{t+1:t+\Delta t-1}|\calE_{:t})\\ 
	&= \sum_{\calE_{t+1:t+\Delta t-1}} \prob(\calE_{t+\Delta t}|\calE_{:t+\Delta t-1})  \cdots  \prob(\calE_{t+1}|\calE_{:t})\\ 
	&= \mathbb{E}_{\calE_{t+1:t+\Delta t-1}|\calE_{:t}} [\prob(\calE_{t+\Delta t}|\calE_{:t+\Delta t-1})] \\
	&\simeq \prob(\calE_{t+\Delta t}|\hat{\calE}_{t+1:t+\Delta t-1},\calE_{:t}).
\end{align*}
\normalsize

Intuitively, one starts with computing {\small $\prob(\calE_{t+1}|\calE_{:t})$}, and drawing a sample $\hat{\calE}_{t+1}$ from the conditional distribution. With this sample, one can further compute {\small $\prob(\calE_{t+2}|\hat{\calE}_{t+1},\calE_{:t})$}. By iteratively computing the conditional distribution for $\calE_{t'}$ and drawing a sample from it, one can eventually estimate {\small $\prob (\calE_{t+\Delta t}|\calE_{:t})$} as {\small $\prob(\calE_{t+\Delta t}|\hat{\calE}_{t+1:t+\Delta t-1},\calE_{:t})$}. 
Although we can improve the estimation by drawing multiple graph samples at each step, \method already performs very well with a single sample, and thus we only draw one sample graph at each step for better efficiency. 
Based on the estimation of the conditional distribution, we can further predict events that are likely to form in the future. 
We summarize the detailed inference algorithm in Algorithm~\ref{alg:renet}; we first sample $M$ number of $\bs$ (line~\ref{alg:samples}) and pick top-$k$ triples (line~\ref{alg:sampletriple}). Then we build a graph at time $t'$ (line~\ref{alg:graph}) to generate a graph. The time complexity of the algorithm is described in Section~\ref{sec:complexity} of appendix.




